\documentclass[11pt,a4paper]{article}

\usepackage[margin=2.5cm]{geometry}
\usepackage{enumitem}
\usepackage{hyperref}
\usepackage{parskip}

\hypersetup{colorlinks=true, linkcolor=blue, urlcolor=blue}

\title{\textbf{ForumDB: A PostgreSQL-Powered Community Discussion Platform} \\ \Large \textbf{DBIS Course Project Outline}}
\author{
Abhinav Chowdary Bikkina (23B1018) \\
Guthikonda Abhiram (23B1084) \\
U Sai Likith Jadhav (23B1058) \\ 
Kunsoth Chandrashen Nayak (23B1066)
}
\date{}

\begin{document}

\maketitle
\vspace{-1em}

\section*{Project Overview}

ForumDB is a community discussion platform inspired by Reddit, where users create and join topic-based communities, submit posts, and engage in threaded discussions through nested comments and voting. The design emphasis of this project is on demonstrating how advanced PostgreSQL features can be used to implement core application functionality directly within the database layer.

Rather than treating the database as only a storage layer, this project explores how SQL queries can handle tasks such as feed ranking, hierarchical comment traversal, search, and aggregation. This allows PostgreSQL to serve as the primary computational engine of the platform.

\section*{Core Features}

\begin{itemize}[noitemsep, topsep=4pt]
    \item \textbf{Communities and Posts:} Users can create and join topic-based communities and submit posts within them.

    \item \textbf{Nested Comment Threads:} Comments may reply to other comments, forming hierarchical discussion trees. These structures will be retrieved using recursive SQL queries.

    \item \textbf{Voting System:} Users can upvote or downvote posts and comments. Vote integrity will be enforced using database constraints and triggers to ensure that each user can vote only once per item.

    \item \textbf{Ranked Feeds:} Posts will be displayed using multiple ranking strategies such as \textit{Hot}, \textit{Rising}, and \textit{Controversial}. These rankings will be computed using SQL scoring formulas based on vote counts and post timestamps.

    \item \textbf{Full-Text Search:} PostgreSQL full-text search functionality will allow users to search across post titles and content using \texttt{tsvector} indexing.

    \item \textbf{User Karma:} A simple reputation score will be computed from the voting history of a user's posts and comments.
\end{itemize}

\section*{Database Concepts Demonstrated}

This project aims to demonstrate several important PostgreSQL capabilities:

\begin{itemize}[noitemsep, topsep=4pt]
    \item \textbf{Recursive CTEs:} Used to retrieve arbitrarily deep comment trees through recursive traversal of parent-child relationships.

    \item \textbf{Full-Text Search:} Implemented using PostgreSQL \texttt{tsvector} representations and \texttt{GIN} indexes to support efficient keyword search.

    \item \textbf{Triggers and Constraints:} Used to enforce vote integrity and maintain consistency in post score updates.

    \item \textbf{Window Functions and Materialized Views:} Window functions will be used for feed ranking and karma computation across a user's post and comment history. Community-level statistics such as post counts and activity metrics will be maintained as materialized views.

    \item \textbf{Indexing Strategies:} B-tree indexes will be used for primary keys and foreign keys, while GIN indexes will support full-text search queries.
\end{itemize}

\section*{System Architecture and Implementation Plan}

The system will consist of three main components:

\begin{enumerate}[noitemsep, topsep=4pt]
    \item \textbf{PostgreSQL Database:} Stores all application data and performs ranking, aggregation, and hierarchical queries.

    \item \textbf{Backend Server:} A lightweight backend implemented using Python (Flask) that exposes REST APIs for creating posts, comments, votes, and retrieving feeds.

    \item \textbf{Frontend Interface:} A simple web interface used to demonstrate the platform's functionality, including community browsing, post feeds, comment threads, and search.
\end{enumerate}

The implementation will proceed in stages: designing the relational schema (users, communities, posts, comments, votes), implementing core CRUD operations, implementing recursive comment retrieval, implementing feed ranking queries, adding full-text search support, and finally building a minimal frontend for demonstration.

\section*{Use of Generative AI Tools}

Generative AI tools may be used to assist with SQL query drafting, schema design suggestions, and boilerplate code generation. All such usage will be explicitly documented in the final report, including the prompts used and the sections of the project where these tools contributed.

\end{document}